\documentclass[11pt]{article}
\usepackage[letterpaper,margin=1in]{geometry}
\usepackage[T1]{fontenc}
\usepackage[utf8]{inputenc}
\usepackage{lmodern,microtype}
\usepackage{amsmath,amssymb,mathtools,bm,slashed,mathrsfs}
\usepackage{xcolor,enumitem,booktabs,array,needspace}
\usepackage[most]{tcolorbox}
\usepackage{fancyhdr,titlesec,parskip,hyperref}
\definecolor{qftblue}{HTML}{245A8D}
\definecolor{qftgreen}{HTML}{2F6B4F}
\definecolor{qftgray}{HTML}{F4F6F8}
\hypersetup{colorlinks=true,linkcolor=qftblue,urlcolor=qftblue,pdfauthor={Nima Lashkari}}
\pagestyle{fancy}\fancyhf{}
\lhead{PHYS 662: Quantum Field Theory I}\cfoot{\thepage}
\setlength{\headheight}{14pt}
\titleformat{\section}{\Large\bfseries\color{qftblue}}{\thesection.}{0.6em}{}
\titleformat{\subsection}{\large\bfseries\color{qftgreen}}{\thesubsection}{0.6em}{}
\setlist[enumerate]{leftmargin=2em,itemsep=0.25em,topsep=0.35em}
\setlist[itemize]{leftmargin=2em,itemsep=0.2em,topsep=0.35em}
\newcommand{\dd}{\mathrm d}
\newcommand{\ii}{\mathrm i}
\newcommand{\eps}{\varepsilon}
\newcommand{\tr}{\operatorname{tr}}
\newcommand{\Tr}{\operatorname{Tr}}
\newcommand{\diag}{\operatorname{diag}}
\newcommand{\Lag}{\mathcal L}
\newcommand{\D}{\mathcal D}
\newcommand{\vac}{\lvert0\rangle}
\newcommand{\vbra}{\langle0\rvert}
\newcommand{\Tord}{\mathrm T}
\newcommand{\ket}[1]{\lvert#1\rangle}
\newcommand{\bra}[1]{\langle#1\rvert}
\newtcolorbox{keyidea}[1][]{enhanced,breakable,colback=qftgray,colframe=qftblue,
 boxrule=0.7pt,arc=1.5mm,title={#1},fonttitle=\bfseries,coltitle=white}
\newtcolorbox{remarkbox}[1][]{enhanced,breakable,colback=white,colframe=qftgreen,
 boxrule=0.6pt,arc=1.5mm,title={#1},fonttitle=\bfseries,coltitle=white}
\newcommand{\intp}{\int\!\frac{\dd^4p}{(2\pi)^4}}
\newcommand{\MSbar}{\overline{\mathrm{MS}}}
\newcommand{\module}[1]{\clearpage\section{#1}}
\newcommand{\braket}[2]{\langle#1\mid#2\rangle}
\newcommand{\dPi}{\dd\Pi}
\newcommand{\cM}{\mathcal M}
\newcommand{\cH}{\mathcal H}
\newcommand{\cS}{\mathcal S}
\newcommand{\Zp}{Z_{\rm pole}}
\newcommand{\Kallen}{\mathcal K}
\rhead{Homework Set 9}
\hypersetup{pdftitle={PHYS 662 Homework Set 9: Quantum Electrodynamics}}
\newlist{parts}{enumerate}{1}
\setlist[parts]{label=\textbf{(\alph*)},leftmargin=2.3em,itemsep=0.55em,topsep=0.4em}
\newcommand{\problem}[2]{\clearpage\section*{\color{qftblue}Problem #1: #2}}
\newcommand{\connection}[1]{{\small\textit{Lecture connection: #1}\par}\medskip}
\begin{document}
\begin{center}
{\scshape\Large PHYS 662 --- Quantum Field Theory I\par}
\vspace{0.3em}{\Huge\bfseries\color{qftblue}Homework Set 9\par}
\vspace{0.3em}{\large Quantum Electrodynamics: Photons, Amplitudes, and Loops\par}
\vspace{0.3em}{\large Lectures 26--29 \quad Fall 2026\par}
\end{center}\medskip
\textbf{Name:} \rule{0.43\textwidth}{0.4pt}\hfill
\textbf{Due:} \rule{0.24\textwidth}{0.4pt}

\textbf{Instructions.} Complete all four problems, worth 25 points each.
Show the steps behind gauge cancellations and normalization factors.
Distinguish ultraviolet subtraction from infrared regulation, and state
where on-shell conditions are used. You may use a supplied intermediate
result even if you have not derived it. No full evaluation of the
off-shell vertex or long Compton gamma trace is required.

\begin{keyidea}[Common conventions]
Use \(\hbar=c=1\), \(\eta=\diag(1,-1,-1,-1)\),
\(\alpha=e^2/(4\pi)\), and signed charge \(e<0\) for an electron.
\[
 \Lag_{\rm int}=-e\bar\psi\gamma^\mu\psi A_\mu,\qquad
 S_F(p)=\frac{\ii(\slashed p+m)}{p^2-m^2+\ii0},\qquad
 D_{\mu\nu}^{(1)}(k)=\frac{-\ii\eta_{\mu\nu}}{k^2+\ii0}.
\]
Feynman rules compute \(\ii\cM\), where
\(\langle f|S-1|i\rangle_c=\ii(2\pi)^4\delta^{(4)}(P_f-P_i)\cM\).
Use covariant one-particle normalization and
\(\dd\Pi_p=\dd^3p/[(2\pi)^3 2E_p]\). The loop conventions are
\[
 d=4-2\eps,\quad \frac1{\bar\eps}=\frac1\eps-\gamma_E+\ln4\pi,
 \quad a=\frac{\alpha}{4\pi},
\]
\[
 \ii\Pi^{\mu\nu}=\ii(q^2\eta^{\mu\nu}-q^\mu q^\nu)\Pi(q^2),\qquad
 S_R=\frac{\ii}{\slashed p-m-\Sigma_R(p)}.
\]
The photon Dyson denominator is \(1-\Pi_R\). The vertex is
\(-\ii e\Gamma_R^\mu\), with photon momentum \(q=p'-p\).
For on-shell spinors, \(\bar u u=2M\),
\(\sum_su_s\bar u_s=\slashed p+M\), and
\(\sigma^{\mu\nu}=\ii[\gamma^\mu,\gamma^\nu]/2\).
\end{keyidea}

\textbf{Coverage.} Problem 1 connects Maxwell constraints to physical
photon states. Problem 2 turns a gauge-invariant QED amplitude into an
observable. Problem 3 treats the two propagator corrections. Problem 4
connects the vertex identity, magnetic moment, and soft radiation.

\textbf{Useful integrals.}
\[
 \int_0^1x(1-x)\,\dd x=\frac16,\quad
 \int_0^1x^2(1-x)^2\dd x=\frac1{30},\quad
 \int_0^1x(1-x)\ln[x(1-x)]\dd x=-\frac5{18}.
\]
All charged-particle masses are nonzero. Where a pole residue or on-shell
vertex requires an infrared regulator, keep it until the indicated
physical combination is formed.

\problem{1}{Why a photon has two polarizations}
\connection{Lecture 26, Sections 1--8.}
Consider Maxwell theory with \(\Lag=-F^2/4-j^\mu A_\mu\),
\(j^\mu=(\rho,\bm J)\), and \(\partial_\mu j^\mu=0\). Fields and gauge
transformations vanish suitably at spatial infinity.
\begin{parts}
\item \textbf{(4 points)} Compute the momentum conjugate to \(A_0\) and
derive Gauss's law by varying \(A_0\). Explain why \(A_0\) is not an
independent oscillator. What equation would be lost if one set \(A_0=0\)
before variation and imposed no further constraint?

\item \textbf{(5 points)} In Coulomb gauge,
\(\bm\nabla\cdot\bm A=0\), solve for \(A^0\). Substitute it into the
action to derive the Coulomb term in the reduced Lagrangian, including
its sign and factor \(1/2\). Check that two static charges have interaction
energy \(q_1q_2/(4\pi r)\) after omitting their separate self-energies.

\item \textbf{(4 points)} With
\(P_T^{ij}=\delta^{ij}-k^ik^j/|\bm k|^2\), show that the two-polarization
mode expansion and
\([a_\lambda(\bm k),a_{\lambda'}^\dagger(\bm q)]
=2|\bm k|(2\pi)^3\delta_{\lambda\lambda'}\delta^{(3)}(\bm k-\bm q)\)
imply \([A_T^i(\bm x),\dot A_T^j(\bm y)]=\ii\delta_T^{ij}(\bm x-\bm y)\).
Why would an unrestricted spatial delta function be inconsistent?

\item \textbf{(4 points)} For a null \(k\), count the independent
polarizations satisfying \(k\cdot\epsilon=0\), with
\(\epsilon\sim\epsilon+ck\). Using a reference vector \(n\) with
\(n\cdot k\ne0\), verify that
\[
 P^{\mu\nu}=-\eta^{\mu\nu}
 +\frac{k^\mu n^\nu+n^\mu k^\nu}{k\cdot n}
 -\frac{n^2k^\mu k^\nu}{(k\cdot n)^2}
\]
is transverse to both \(k\) and \(n\), and reduces to the spatial
polarization sum for \(n=(1,\bm0)\).

\item \textbf{(4 points)} Add \(-(\partial\cdot A)^2/(2\xi)\), invert
the quadratic operator using \(P_T=\eta-kk/k^2\) and \(P_L=kk/k^2\),
and derive \(D_{\mu\nu}^{(\xi)}\). Show that exchange between conserved
currents is independent of \(\xi\). Why is \(P_T\) here not an on-shell
polarization sum?

\item \textbf{(4 points)} In covariant quantization, explain why
\((\partial\cdot A)^{(+)}|\mathrm{phys}\rangle=0\) is imposed instead
of a strong operator condition. Identify a null timelike--longitudinal
one-photon combination for momentum along \(z\), and explain why the
physical Hilbert space requires a quotient by null states.
\end{parts}

\problem{2}{Compton scattering: gauge invariance and recoil}
\connection{Lecture 27, Sections 2, 4, and 5.}
Study \(e^-(p)+\gamma(k)\to e^-(p')+\gamma(k')\), with
\(p^2=p'^2=m^2\), \(k^2=k'^2=0\), and \(p+k=p'+k'\).
\begin{parts}
\item \textbf{(5 points)} Draw the two tree diagrams and derive
\[
 \cM=-e^2\bar u(p')\left[
 \slashed\epsilon'^{\,*}\frac{\slashed p+\slashed k+m}{2p\cdot k}
 \slashed\epsilon+
 \slashed\epsilon\frac{\slashed p-\slashed k'+m}{-2p\cdot k'}
 \slashed\epsilon'^{\,*}\right]u(p).
\]
Explain the relative sign and why minimal Dirac QED has no extra
contact diagram with two photons on the same vertex.

\item \textbf{(4 points)} Replace \(\epsilon\) by \(k\) and demonstrate
explicit cancellation between the two diagrams using the external
Dirac equations. Explain how to obtain the outgoing-photon Ward identity.
Does either individual diagram obey these identities?

\item \textbf{(4 points)} Write the spin-averaged squared amplitude as
a Dirac trace and specify the initial-state average. Explain why the
physical photon sums may now be replaced by \(-\eta_{\mu\nu}\).
Would this replacement be justified on one diagram considered alone?
You do not need to evaluate the long trace.

\item \textbf{(5 points)} In the electron rest frame set
\(p=(m,\bm0)\), \(k^0=\omega\), \(k'^0=\omega'\), with angle
\(\theta\) between \(\bm k\) and \(\bm k'\). Derive
\[
 \omega'=\frac{\omega}{1+(\omega/m)(1-\cos\theta)}.
\]
Starting from invariant two-body phase space and incoming flux
\(4m\omega\), show that
\[
 \frac{\dd\sigma}{\dd\Omega}
 =\frac{\overline{|\cM|^2}}{64\pi^2m^2}
 \left(\frac{\omega'}{\omega}\right)^2.
\]
Display the energy-delta-function Jacobian rather than quoting a
center-of-momentum formula in the laboratory frame.

\item \textbf{(4 points)} You may use
\(\overline{|\cM|^2}=2e^4[\omega/\omega'+\omega'/\omega-\sin^2\theta]\).
Obtain the unpolarized differential cross section and its limit for
\(\omega\ll m\). Integrate that limit over the full solid angle to find
the total Thomson cross section.

\item \textbf{(3 points)} At fixed nonzero scattering angle, determine
the leading dependence of \(\dd\sigma/\dd\Omega\) on \(\omega\) for
\(\omega\gg m\). Explain why this approximation is not uniform in the
forward direction and cannot be integrated over all angles without care.
\end{parts}

\problem{3}{What the two propagator corrections measure}
\connection{Lecture 28, Sections 1--5 and 7--9.}
For one Dirac species, the zero-momentum-subtracted photon loop is
\[
 \widehat\Pi(q^2)=\frac{2\alpha}{\pi}\int_0^1\dd x\,x(1-x)
 \ln\frac{m^2-x(1-x)q^2-\ii0}{m^2}.
\]
\begin{parts}
\item \textbf{(4 points)} Starting with the trace integral in Lecture 28,
contract \(q_\mu\) into the photon 1PI insertion and show that it vanishes
by a shift of regulated loop momentum. State the property the regulator
must have. Explain why this transversality does not imply
\(q^\mu D_{\mu\nu}=0\) in Feynman gauge.

\item \textbf{(5 points)} Sum the transverse Dyson series, then derive
the leading small-\(Q\) and large-\(Q\) limits of \(\widehat\Pi(-Q^2)\).
Use them to discuss screening, decoupling of a heavy fermion, and the
sign of the one-loop running of
\(\alpha_{\rm eff}(Q^2)=\alpha(0)/[1-\widehat\Pi(-Q^2)]\).
Keep the constant accompanying the large logarithm.

\item \textbf{(4 points)} Determine the first branch point in timelike
\(q^2=s\). Find the parameter interval on which the logarithm has an
imaginary part and calculate \(\operatorname{Im}\widehat\Pi(s+\ii0)\).
Relate its mass-dependent factor to the tree-level pair-production
cross section in Lecture 27. Explain the sign in our tensor convention.

\item \textbf{(4 points)} For the electron self-energy take the supplied
ultraviolet result
\[
 \Sigma_{\rm UV}^{(\xi)}(p)
 =\frac{a}{\bar\eps}[-\xi\slashed p+(3+\xi)m].
\]
With \(\Sigma_R=\Sigma_{\rm loop}-\delta_2(\slashed p-m)+m\delta Z_m\),
find the \(\MSbar\) constants \(\delta_2\) and \(\delta Z_m\).
Which depends on the covariant gauge parameter, and why is a vanishing
field ultraviolet pole insufficient to make the self-energy vanish?

\item \textbf{(4 points)} The finite on-shell mass shift in Feynman gauge
can be written
\[
 \frac{M-m(\mu)}{m(\mu)}
 =a\int_0^1\dd x\left[(2+2x)
 \left(\ln\frac{\mu^2}{m(\mu)^2}-2\ln x\right)-2x\right].
\]
Evaluate it. Determine \(\mu\,\dd m/\dd\mu\) to order \(\alpha\) by
requiring the pole mass \(M\) to be independent of \(\mu\).

\item \textbf{(4 points)} For
\(\Sigma_R=A(p^2)\slashed p+B(p^2)m\), derive the mass condition and
regulated pole residue by differentiating the positive-energy inverse
propagator at rest. Explain why the on-shell residue can be infrared
singular even though the one-loop mass shift is finite, and why it is
not generally identical to the \(\MSbar\) field constant.
\end{parts}

\problem{4}{The vertex identity and a finite magnetic moment}
\connection{Lecture 29, Sections 1--8.}
Define \(\mathcal S_R^{-1}(p)=\slashed p-m-\Sigma_R(p)\) and
\(\Gamma_R^\mu=\gamma^\mu+\Lambda^\mu+\delta_1\gamma^\mu\).
\begin{parts}
\item \textbf{(4 points)} Explain the origin of contact terms when a
current divergence acts on the time-ordered product of \(j^\mu(x)\), \(\psi(y)\), and \(\bar\psi(z)\).
Use the resulting Ward--Takahashi identity
\[q_\mu\Gamma_R^\mu=\mathcal S_R^{-1}(p')-\mathcal S_R^{-1}(p)\]
to derive the zero-transfer derivative identity. Why is the off-shell
contraction not generally zero?

\item \textbf{(4 points)} Use Problem 3's ultraviolet self-energy to
predict the vertex pole in a general covariant gauge. Determine
\(\delta_1\), verify \(Z_1=Z_2\) at this order, and explain how this
reduces \(e_0=\mu^\eps eZ_1/(Z_2\sqrt{Z_3})\). Does the absence of the
one-loop vertex ultraviolet pole in Landau gauge imply \(F_2=0\)?

\item \textbf{(4 points)} Derive the Gordon identity and apply it to
\[
 \bar u'\left[\gamma^\mu F_1(q^2)
 +\frac{\ii\sigma^{\mu\nu}q_\nu}{2M}F_2(q^2)\right]u.
\]
Show that \(g=2[F_1(0)+F_2(0)]\) and hence
\(a_e=F_2(0)\) when the physical charge is defined by \(F_1(0)=1\).
Why does the Ward contraction leave the Pauli term unconstrained?

\item \textbf{(5 points)} Combine the three denominators of the vertex
loop with parameters \(x,y\) on the electrons and \(z\) on the photon.
For \(x+y+z=1\), \(p^2=p'^2=M^2\), show that shifting
\(\ell=k-xp-yp'\) gives
\(\Delta=M^2(1-z)^2-xyq^2\). You may use the finite projection
\[
 F_2(q^2)=\frac{\alpha}{\pi}\int_0^1\dd x\,\dd y\,\dd z\,
 \delta(1-x-y-z)\frac{M^2z(1-z)}{\Delta-\ii0}.
\]
Evaluate \(F_2(0)\), explicitly accounting for the parameter-triangle
measure, and give the leading correction to \(g\).

\item \textbf{(4 points)} As an infrared check, replace
\(\Delta\) by \(\Delta+z\lambda^2\) in this finite integral. Reduce
\(F_2(0;\lambda)\) to one dimension and justify the limit
\(\lambda\to0\), including the endpoint \(z=1\).
State the dimension and coefficient of the local Pauli operator that
reproduces this response at small transfer.

\item \textbf{(4 points)} For massive charged external particles, the
soft-emission factor is
\(\sum_i\eta_i e_i p_i\cdot\epsilon^*/(p_i\cdot k)\), with
\(\eta_i=+1\) outgoing and \(-1\) incoming.
Show that charge conservation makes it gauge invariant and that its
phase-space integral behaves as \(\int\dd\omega/\omega\).
State which measured combination cancels the corresponding virtual
infrared singularity and which detector-dependent scale remains.
Why does ultraviolet renormalization alone not perform this cancellation?
\end{parts}
\end{document}
